
\documentclass[11pt]{article}
\usepackage[utf8]{inputenc}
\usepackage{amsmath, amssymb}

\title{Méthodes fondamentales – Spécialité Mathématiques}
\date{}

\begin{document}

\maketitle

\section{Raisonnement par récurrence}
\textbf{Méthode :}
\begin{itemize}
\item Initialisation : vérifier $P(n_0)$
\item Hérédité : supposer $P(k)$ vraie et montrer $P(k+1)$
\item Conclusion
\end{itemize}

\textbf{Exemple :} Montrer que $1+2+\cdots+n=\frac{n(n+1)}{2}$.

\section{Monotonie d’une suite}
\textbf{Méthode :} Étudier $u_{n+1}-u_n$.

\textbf{Exemple :} $u_n = n^2$ \\
$u_{n+1}-u_n = 2n+1 > 0$ donc croissante.

\section{Limite d’une suite}
\textbf{Méthode :} Identifier le terme dominant.

\textbf{Exemple :}
\[
\frac{3n^2+1}{5n^2-4} \to \frac{3}{5}
\]

\section{Croissances comparées}
\textbf{Hiérarchie :} $\ln n \ll n^k \ll q^n \ll n!$

\textbf{Exemple :}
\[
\frac{n^5}{2^n} \to 0
\]

\section{Suites récurrentes}
\textbf{Méthode :} Chercher le point fixe.

\textbf{Exemple :} $u_{n+1} = \frac{1}{2}u_n +3$ \\
$\ell = 6$

\section{Dérivation}
\textbf{Exemple :}
\[
f(x)=\ln(x^2+1), \quad f'(x)=\frac{2x}{x^2+1}
\]

\section{Convexité}
\textbf{Méthode :} Étudier $f''$.

\textbf{Exemple :} $f(x)=x^2 \Rightarrow f''(x)=2>0$ donc convexe.

\section{Continuité – TVI}
\textbf{Méthode :} Continuité + changement de signe.

\textbf{Exemple :} $f(x)=x^3-2$ sur $[1,2]$.

\section{Loi binomiale}
\[
P(X=k)=\binom{n}{k}p^k(1-p)^{n-k}
\]

\section{Équations différentielles}
\[
y'=ay \Rightarrow y=Ce^{ax}
\]

\section{Intégration}
\[
\int_0^1 x\,dx = \frac{1}{2}
\]

\end{document}
